\begin{abstract}
Sharpness-Aware Minimization (SAM) consistently improves generalization across deep-learning settings, including under distribution shift, yet the mechanism behind its robustness to corrupted inputs remains unclear. Existing explanations attribute SAM's benefits to convergence to flatter loss-landscape minima, but several recent results --- notably the finding that perturbing only normalization layers improves generalization while \emph{increasing} measured sharpness --- suggest flatness alone does not fully account for the observed gains. In this paper, we investigate whether SAM improves out-of-distribution (OOD) generalization primarily through loss-landscape sharpness reduction or through learning more stable internal representations that persist under input corruptions. We hypothesise that SAM's perturbation objective encourages representation robustness beyond what sharpness measures capture. To test this, we compare SGD, SAM ($\rho = 0.05$), and Stochastic Weight Averaging (SWA, a flat-minima-oriented control that reaches wide basins via late-epoch trajectory averaging rather than adversarial perturbation) under matched training budgets on CIFAR-10 and CIFAR-100, evaluating in-distribution accuracy, severity-averaged corruption accuracy on CIFAR-10-C and CIFAR-100-C, expected calibration error, top Hessian eigenvalue sharpness, and clean-to-corrupted paired-image CKA stability at the penultimate layer. Rather than asking only whether SAM improves accuracy, we ask whether its OOD gain over SGD is fully explained by sharpness reduction or whether it requires an additional representation-stability explanation. This is the first controlled comparison of SAM and SWA on corruption-robustness benchmarks that jointly measures sharpness and representation stability, providing a decomposition that the existing SAM-for-OOD literature has not previously enabled.
\end{abstract}
